\medskip\noindent\SGARefTarget{sga3:xi:remark:1:7}\textbf{注 1.7。} --- 在 \SGARefLink{sga3:xi:remark:1:6}{1.6} 的假设下，设给定“\(F\) 的一个取值于域 \(k\) 的点”，即元素 \(x\in F(\operatorname{Spec}k)\)，或等价地，态射 \[
  x:\operatorname{Spec}k\longrightarrow F.
\]。若前一定义中的条件在 \(S'\) 为局部概形、且上图 \(\mathrm{(Q)}\) 中的态射 \(S'_0\to F\) 以下述意义“与标记点相容”时成立，则同样称 \(f\) \emph{在 \(x\) 处形式平滑}（分别为\emph{形式非分歧}、\emph{形式平展}）。以 \(k'\) 表示 \(S'\) 和 \(S'_0\) 在闭点处的剩余域。则图

\SGARefTarget{sga3:xi:diagram:1:7:marked-point-compatibility} \[
\begin{tikzcd}[column sep=huge,row sep=large]
  \operatorname{Spec}(k') \arrow[d,"j"'] &
    \operatorname{Spec}(k) \arrow[d,"x"] \\
  S'_0 \arrow[r,"u_0"] & F
\end{tikzcd}
\]，其中 \(j:\operatorname{Spec}(k')\to S'_0\) 为典范浸入，必须能够补成交换图

\SGARefTarget{sga3:xi:diagram:1:7:common-field-extension} \[
\begin{tikzcd}[column sep=huge,row sep=large]
  & \operatorname{Spec}(k'') \arrow[dl] \arrow[dr] & \\
  \operatorname{Spec}(k') \arrow[d] &&
    \operatorname{Spec}(k) \arrow[d] \\
  S'_0 \arrow[rr] && F ,
\end{tikzcd}
\]，其中 \(k''\) 为域。\footnote{原文称 \(k''\) 为“一个域的谱”；所显示的对象 \(\operatorname{Spec}(k'')\) 表明 \(k''\) 本身是域。}

当 \(F\) 和 \(G\) 可由 \(S\)-预概形 \(X\) 和 \(Y\) 表示，且 \(S\)-态射 \(f:X\to Y\) 局部有限呈示时，由 \SGARefLink{sga3:xi:remark:1:5}{1.5}，此条件恰是说 \(f:X\to Y\) 在 \(\operatorname{Spec}(k)\) 经 \(x:\operatorname{Spec}(k)\to X\) 所映到的 \(X\) 的点处平滑。

\medskip\noindent\SGARefTarget{sga3:xi:remark:1:8}\textbf{注 1.8。} --- 若前一定义中的条件在只考虑局部 Artin 概形 \(S'\) 后成立，则称 \(F\to G\) \emph{在 \(x\) 处无穷小平滑}（分别为\emph{无穷小非分歧}、\emph{无穷小平展}）。若 \(F\to G\) 在每个点 \(x\) 处均具有相应性质，则称它无穷小平滑（分别为无穷小非分歧、无穷小平展）；等价地，\SGARefLink{sga3:xi:remark:1:6}{1.6} 所考虑的条件在 \(S'\) 为局部 Artin 概形时总成立。

上述概念的这一变体在技术上很有用。由于它较弱，往往更容易验证，却仍经常足以推出较强的条件；例如，当 \(F\to G\) 为局部有限呈示态射，且 \(G\) 可由局部 Noether 预概形表示时就是如此。


